% =============================================================================
% LaTeX-VORLAGE: Projektdokumentation Abschlusspruefung FIAE (Teil 1 oder Teil 2)
%
% Wiederverwendbare, projektunabhaengige Vorlage. Orientiert an:
%  - Checkliste_zur_Projektdokumentation.docx (IHK)
%  - Beispieldokumentation Berg (2025/26) -- gleiche Kapitel-/Anhangstruktur
%  - latex-vorlage-fiae-master (Stefan Macke) -- Grundaufbau
%
% Kompilieren:  latexmk -pdf main.tex   (oder: pdflatex main.tex, zweimal)
%
% VOR DER NUTZUNG:
%  1. Alle [PLATZHALTER: ...] in diesem Verzeichnis suchen & ersetzen
%     (z.B. per Editor-Suche ueber alle Dateien).
%  2. kapitel/*.tex mit echtem Inhalt fuellen (Ueberschriften/Struktur stehen).
%  3. anhang/*.tex: jede Anlage entweder mit echtem Inhalt fuellen oder
%     (falls im eigenen Projekt nicht zutreffend) Datei samt \input-Zeile
%     in anhang.tex entfernen.
%  4. Seitenumfang pruefen: Hauptteil max. 15 DIN-A4-Seiten, alle Anlagen
%     zusammen max. 20 Seiten (siehe README.md).
% =============================================================================
\documentclass[
  ngerman,
  11pt,
  a4paper,
  oneside,
  parskip=half*,
  headings=normal,
  listof=totoc,
  bibliography=totoc,
]{scrreprt}

% --- Sprache & Encoding -------------------------------------------------------
\usepackage[utf8]{inputenc}
\usepackage[T1]{fontenc}
\usepackage[ngerman]{babel}
\usepackage{csquotes}

% --- Schrift: Arial oder Tahoma gefordert (11pt, siehe README). Helvetica-
%     Metrik-Klon als freie, ueberall verfuegbare Alternative (Overleaf/TeX Live
%     ohne Font-Lizenz nutzbar). Bei lokal installiertem Arial/Tahoma stattdessen
%     mit lualatex + fontspec (\setmainfont{Arial}) kompilieren. ------------------
\usepackage{helvet}
\renewcommand{\familydefault}{\sfdefault}
\usepackage{microtype}

% --- Satzspiegel nach DIN 5008 -------------------------------------------------
\usepackage[left=3.5cm,right=2.5cm,top=2.5cm,bottom=2.5cm,headsep=1cm,footskip=1.2cm]{geometry}

% --- Zeilenabstand 1,3 (Vorgabe) ------------------------------------------------
\usepackage{setspace}
\setstretch{1.3}

% --- Grafiken & Diagramme (TikZ) -----------------------------------------------
\usepackage{graphicx}
\usepackage{xcolor}
\usepackage{tikz}
\usetikzlibrary{arrows.meta, positioning, shapes.geometric, fit, calc}

% --- Tabellen -------------------------------------------------------------------
\usepackage{booktabs}
\usepackage{array}
\usepackage{caption}
\captionsetup{font=small, labelfont=bf}

% --- Code-/Listingsverzeichnis ---------------------------------------------------
\usepackage{listings}
\lstdefinestyle{code}{
  basicstyle=\ttfamily\small,
  keywordstyle=\bfseries,
  commentstyle=\itshape\color{gray!60!black},
  stringstyle=\color{gray!40!black},
  showstringspaces=false,
  breaklines=true,
  frame=single,
  framerule=0.3pt,
  numbers=left,
  numberstyle=\tiny\color{gray},
  xleftmargin=1.5em,
}
\lstset{style=code}

% --- Listen ----------------------------------------------------------------------
\usepackage{enumitem}

% --- Anlagen als PDF einbinden (persoenliche Erklaerung, Protokoll, Formulare) --
\usepackage{pdfpages}

% --- Hurenkinder/Schusterjungen unterdruecken -----------------------------------
\widowpenalty=10000
\clubpenalty=10000

% --- Verlinkung (immer als letztes Paket laden) ---------------------------------
\usepackage[
  colorlinks=true,
  linkcolor=black,
  citecolor=black,
  urlcolor=black,
  bookmarksnumbered=true,
]{hyperref}

% --- Metadaten -------------------------------------------------------------------
\newcommand{\dokTitel}{[PLATZHALTER: Projekttitel]}
\newcommand{\dokUntertitel}{Projektdokumentation zur Abschlusspr\"ufung Teil [PLATZHALTER: 1 oder 2] -- Fachinformatiker/-in Anwendungsentwicklung}
\newcommand{\dokPruefling}{[PLATZHALTER: Name des Pr\"uflings]}
\newcommand{\dokKennnummer}{[PLATZHALTER: Pr\"ufungs-/Kennnummer]}
\newcommand{\dokBetrieb}{[PLATZHALTER: Ausbildungsbetrieb]}
\newcommand{\dokZeitraum}{[PLATZHALTER: TT.MM.JJJJ] -- [PLATZHALTER: TT.MM.JJJJ]}
\newcommand{\dokAbgabedatum}{[PLATZHALTER: TT.MM.JJJJ]}

\hypersetup{
  pdftitle={\dokTitel},
  pdfauthor={\dokPruefling},
  pdflang={de-DE},
}

\begin{document}

% =============================================================================
% VORTEXT (ungezaehlt / kein Seitenstil)
% =============================================================================
\pagestyle{empty}
% Deckblatt -------------------------------------------------------------------
\begin{titlepage}
  \centering
  \vspace*{1.5cm}

  {\large \dokBetrieb\par}
  \vspace{0.3cm}
  {\small [PLATZHALTER: Firmenlogo einf\"ugen, falls vom Ausbildungsbetrieb gew\"unscht]\par}

  \vspace{2.5cm}
  {\Large\bfseries Projektdokumentation\par}
  \vspace{0.3cm}
  {\large zur Abschlusspr\"ufung Teil [PLATZHALTER: 1 oder 2]\par}
  {\large Fachinformatiker/-in Anwendungsentwicklung\par}

  \vspace{2cm}
  {\LARGE\bfseries \dokTitel \par}

  \vspace{2.5cm}
  \begin{tabular}{ll}
    Pr\"ufling:            & \dokPruefling \\[0.3em]
    Kennnummer:            & \dokKennnummer \\[0.3em]
    Ausbildungsbetrieb:    & \dokBetrieb \\[0.3em]
    Bearbeitungszeitraum:  & \dokZeitraum \\[0.3em]
    Ausbildende/r:         & [PLATZHALTER: Name Ausbilder/in] \\[0.3em]
    Abgabedatum:           & \dokAbgabedatum \\
  \end{tabular}

  \vfill
  {\small
    [PLATZHALTER: Falls Teile der Dokumentation vertrauliche oder
    datenschutzrelevante Angaben enthalten (Kundendaten, interne Systeme,
    Kennzahlen etc.), hier auf den Sperrvermerk auf der Folgeseite verweisen
    und die betroffenen Stellen im Text entsprechend kennzeichnen. Falls kein
    Sperrvermerk noetig ist, diesen Absatz sowie sperrvermerk.tex entfernen.]
    \par}
\end{titlepage}

% Sperrvermerk ------------------------------------------------------------------
% Nur noetig, wenn die Dokumentation vertrauliche/datenschutzrelevante Daten
% enthaelt (siehe Vorgabe: "Vertrauliche oder datenschutzrelevante Daten sind
% als solche zu kennzeichnen"). Mit Ausbildungsbetrieb und IHK abstimmen. Falls
% nicht benoetigt: Datei leeren oder % Sperrvermerk ------------------------------------------------------------------
% Nur noetig, wenn die Dokumentation vertrauliche/datenschutzrelevante Daten
% enthaelt (siehe Vorgabe: "Vertrauliche oder datenschutzrelevante Daten sind
% als solche zu kennzeichnen"). Mit Ausbildungsbetrieb und IHK abstimmen. Falls
% nicht benoetigt: Datei leeren oder % Sperrvermerk ------------------------------------------------------------------
% Nur noetig, wenn die Dokumentation vertrauliche/datenschutzrelevante Daten
% enthaelt (siehe Vorgabe: "Vertrauliche oder datenschutzrelevante Daten sind
% als solche zu kennzeichnen"). Mit Ausbildungsbetrieb und IHK abstimmen. Falls
% nicht benoetigt: Datei leeren oder \input{sperrvermerk} in main.tex entfernen.
\thispagestyle{empty}
\vspace*{4cm}

\begin{center}
  \bfseries Sperrvermerk
\end{center}

\vspace{1cm}

\noindent
[PLATZHALTER: Standardformulierung, z.B. "Die vorliegende
Projektdokumentation enth\"alt vertrauliche Informationen von \dokBetrieb,
insbesondere zu [PLATZHALTER: betroffene Bereiche, z.B. interne Systeme,
Kundendaten, Kennzahlen]. Eine Weitergabe oder Ver\"offentlichung des Inhalts,
auch auszugsweise, ist ohne ausdr\"uckliche schriftliche Zustimmung von
\dokBetrieb nicht gestattet. Ausgenommen hiervon ist die Einsichtnahme durch
die zust\"andigen Pr\"ufungsausschuss- und IHK-Mitglieder im Rahmen der
Abschlusspr\"ufung."]

\vspace{2cm}
\noindent
[PLATZHALTER: Ort, Datum, Unterschrift Ausbildungsbetrieb -- vor Abgabe
erg\"anzen bzw. diese Seite entfernen, falls kein Sperrvermerk gefordert wird.]

\clearpage
 in main.tex entfernen.
\thispagestyle{empty}
\vspace*{4cm}

\begin{center}
  \bfseries Sperrvermerk
\end{center}

\vspace{1cm}

\noindent
[PLATZHALTER: Standardformulierung, z.B. "Die vorliegende
Projektdokumentation enth\"alt vertrauliche Informationen von \dokBetrieb,
insbesondere zu [PLATZHALTER: betroffene Bereiche, z.B. interne Systeme,
Kundendaten, Kennzahlen]. Eine Weitergabe oder Ver\"offentlichung des Inhalts,
auch auszugsweise, ist ohne ausdr\"uckliche schriftliche Zustimmung von
\dokBetrieb nicht gestattet. Ausgenommen hiervon ist die Einsichtnahme durch
die zust\"andigen Pr\"ufungsausschuss- und IHK-Mitglieder im Rahmen der
Abschlusspr\"ufung."]

\vspace{2cm}
\noindent
[PLATZHALTER: Ort, Datum, Unterschrift Ausbildungsbetrieb -- vor Abgabe
erg\"anzen bzw. diese Seite entfernen, falls kein Sperrvermerk gefordert wird.]

\clearpage
 in main.tex entfernen.
\thispagestyle{empty}
\vspace*{4cm}

\begin{center}
  \bfseries Sperrvermerk
\end{center}

\vspace{1cm}

\noindent
[PLATZHALTER: Standardformulierung, z.B. "Die vorliegende
Projektdokumentation enth\"alt vertrauliche Informationen von \dokBetrieb,
insbesondere zu [PLATZHALTER: betroffene Bereiche, z.B. interne Systeme,
Kundendaten, Kennzahlen]. Eine Weitergabe oder Ver\"offentlichung des Inhalts,
auch auszugsweise, ist ohne ausdr\"uckliche schriftliche Zustimmung von
\dokBetrieb nicht gestattet. Ausgenommen hiervon ist die Einsichtnahme durch
die zust\"andigen Pr\"ufungsausschuss- und IHK-Mitglieder im Rahmen der
Abschlusspr\"ufung."]

\vspace{2cm}
\noindent
[PLATZHALTER: Ort, Datum, Unterschrift Ausbildungsbetrieb -- vor Abgabe
erg\"anzen bzw. diese Seite entfernen, falls kein Sperrvermerk gefordert wird.]

\clearpage


% =============================================================================
% VERZEICHNISSE (roemische Seitenzahlen)
% =============================================================================
\pagestyle{plain}
\pagenumbering{Roman}

\tableofcontents
\clearpage
\listoffigures
\clearpage
\listoftables
\clearpage
\lstlistoflistings
\clearpage
% Abkuerzungsverzeichnis ---------------------------------------------------------
% Vorgabe: "Nicht allgemeingueltige Abkuerzungen sind zu vermeiden oder
% erforderlichenfalls zu erlaeutern." Nur Abkuerzungen aufnehmen, die im Text
% tatsaechlich verwendet werden -- keine allgemein bekannten (z.B. "z.B.",
% "usw.") auflisten.
\chapter*{Abk\"urzungsverzeichnis}
\addcontentsline{toc}{chapter}{Abk\"urzungsverzeichnis}

\begin{table}[h!]
  \centering
  \begin{tabular}{@{}ll@{}}
    \toprule
    \textbf{Abk\"urzung} & \textbf{Bedeutung} \\
    \midrule
    API   & Application Programming Interface (Programmierschnittstelle) \\
    IHK   & Industrie- und Handelskammer \\
    UI    & User Interface (Benutzeroberfl\"ache) \\
    {[}PLATZHALTER{]} & [PLATZHALTER: weitere projektspezifische Abk\"urzungen erg\"anzen] \\
    \bottomrule
  \end{tabular}
\end{table}

\clearpage

% Roemische Seitenzahl merken, um nach dem Hauptteil dort weiterzuzaehlen
\newcounter{romanWeiter}
\setcounter{romanWeiter}{\value{page}}

% =============================================================================
% HAUPTTEIL (arabische Seitenzahlen) -- max. 15 DIN-A4-Seiten, siehe README.md
% =============================================================================
\pagenumbering{arabic}
\setcounter{page}{1}

\chapter{Einleitung}
\label{ch:einleitung}

% Vorgabe: Der gesamte Projektverlauf von der Ausgangssituation bis zur
% Kundenuebergabe muss transparent dargestellt werden -- keine reine
% Produkt-/Ergebnisbeschreibung. Dieses Kapitel liefert den Rahmen dafuer.

\section{Projektumfeld}
[PLATZHALTER: Ausbildungsbetrieb, Abteilung/Team, fachliches Umfeld, in dem
das Projekt entstanden ist. Wer ist Auftraggeber, wer ist Nutzer/Kunde des
Ergebnisses?]

\section{Projektziel und Projektmotivation}
\label{sec:projektziel}
[PLATZHALTER: Ausgangssituation/Problem, das zum Projekt gefuehrt hat; das
konkrete Projektziel; warum dieses Ziel fuer den Betrieb relevant ist. Bezug
zum genehmigten Projektantrag herstellen (Anlage~\ref{anh:projektantrag}).]

\section{Projektschnittstellen}
[PLATZHALTER: Beteiligte Personen/Rollen (z.B. Ausbilder/in, Fachabteilung,
externe Dienstleister) und technische Schnittstellen zu bestehenden Systemen.]

\section{Projektabgrenzung}
[PLATZHALTER: Was ist explizit NICHT Teil des Projekts? Grenzt das Projekt
gegen Anschlussprojekte oder bestehende Systeme ab.]

\chapter{Anforderungsanalyse}
\label{ch:anforderungsanalyse}

\section{Ist-Analyse}
[PLATZHALTER: Bestehende Situation vor dem Projekt: vorhandene Systeme,
Prozesse, Schwachstellen. Grundlage fuer die Zielsetzung aus
Abschnitt~\ref{sec:projektziel}.]

\section{Lastenheft}
\label{sec:lastenheft}
[PLATZHALTER: Kurzdarstellung der Anforderungen aus Nutzersicht (Was wird
gebraucht und warum?). Vollstaendiges Lastenheft als
Anlage~\ref{anh:lastenheft} beigefuegt -- hier nur Kernanforderungen und
Verweis darauf.]

\section{Soll-Konzept}
[PLATZHALTER: Loesungsidee auf konzeptioneller Ebene, noch ohne technische
Details. Falls mehrere Loesungsalternativen erwogen wurden: hier oder in
Kapitel~\ref{ch:projektplanung} vergleichend darstellen und die Entscheidung
begruenden (Vorgabe: Loesungsalternativen und Entscheidungen sind zu
erlaeutern, nicht nur das Ergebnis).]

\section{Technische Grundlagen}
[PLATZHALTER: Fuer das Verstaendnis noetige Technologien/Konzepte kurz
erlaeutern (nur was fuer Kapitel~\ref{ch:umsetzung} vorausgesetzt wird).
Nicht allgemeingueltige Fachbegriffe/Abkuerzungen hier erklaeren oder ins
Abkuerzungsverzeichnis aufnehmen.]

\section{Pflichtenheft}
\label{sec:pflichtenheft}
[PLATZHALTER: Aus dem Lastenheft abgeleitete technische Umsetzungsvorgaben
(Wie wird es umgesetzt?). Vollstaendiges Pflichtenheft als
Anlage~\ref{anh:pflichtenheft} beigefuegt -- hier nur Kernpunkte und Verweis
darauf.]

\chapter{Projektplanung}
\label{ch:projektplanung}

% Vorgabe: "Mit der Durchfuehrung des Projektes kann erst nach der
% Genehmigung begonnen werden." -- Genehmigungsdatum hier oder in Kap. 7.2
% nachvollziehbar dokumentieren.
[PLATZHALTER: Der Projektantrag wurde vor Beginn der Umsetzung bei der IHK
eingereicht und erst nach dessen Genehmigung (Genehmigungsdatum:
[PLATZHALTER: TT.MM.JJJJ]) begonnen.]

\section{Zeitplanung}
\label{sec:zeitplanung}
[PLATZHALTER: Geplante Phasen und Stunden gemaess Projektantrag kurz
zusammenfassen. Detaillierter Zeitplan/Gantt-Diagramm als
Anlage~\ref{anh:zeitplanung}/\ref{anh:gantt} beigefuegt.]

\section{Ressourcen- und Kostenplanung}
[PLATZHALTER: Benoetigte Ressourcen (Personal, Hardware, Software,
Lizenzen). Kostenaufstellung als Anlage~\ref{anh:kostenplanung} beigefuegt --
hier Kernzahlen und ggf. Amortisationsrechnung/Wirtschaftlichkeitsbetrachtung
zusammenfassen.]

\section{Risikoanalyse}
[PLATZHALTER: Wesentliche Projektrisiken (technisch, organisatorisch,
zeitlich) und Gegenmassnahmen. Vollstaendige Analyse als
Anlage~\ref{anh:risikoanalyse} beigefuegt.]

\section{Stakeholderanalyse}
[PLATZHALTER: Relevante Stakeholder, ihre Interessen/Einfluss auf das
Projekt. Vollstaendige Matrix als Anlage~\ref{anh:stakeholderanalyse}
beigefuegt.]

\chapter{Vorgehensmodell und Methodik}
\label{ch:vorgehensmodell}

\section{Vorgehensmodell}
[PLATZHALTER: Gewaehltes Vorgehensmodell (z.B. Wasserfall, Scrum, Kanban,
individuell) und Begruendung, warum es fuer dieses Projekt passend war.
Details als Anlage~\ref{anh:vorgehensmodell} beigefuegt.]

\section{Programmiersprache, Entwicklungstools und Schnittstellen}
[PLATZHALTER: Eingesetzte Sprachen, Frameworks, Entwicklungsumgebung,
Versionsverwaltung, externe Schnittstellen/APIs -- jeweils kurz begruendet.]

\section{Entwurf von Testfaellen}
[PLATZHALTER: Testfallkonzept: welche Testarten (z.B. Modultest,
Integrationstest, Abnahmetest), nach welchen Kriterien Testfaelle definiert
wurden. Durchfuehrung und Ergebnisse in
Kapitel~\ref{ch:testphase_und_fehlerbehebung} bzw.
Anlage~\ref{anh:testprotokoll}.]

\chapter{Umsetzung}
\label{ch:umsetzung}

% Vorgabe: Abschnitte 5.3-5.5 an die tatsaechlichen Kernmodule des eigenen
% Projekts anpassen (umbenennen, Anzahl variieren) -- nicht wortwoertlich
% uebernehmen. Interessante Quelltext-Ausschnitte hier per \lstinputlisting
% einbinden oder auf Anlage~\ref{anh:quellcode} verweisen; vollstaendiger
% Quelltext gehoert NICHT in den Hauptteil (Seitenlimit).

\section{Datenbank-/Datenmodelldesign}
[PLATZHALTER: Datenmodell, Tabellen-/Entitaetsstruktur. Vollstaendiges
ER-Modell/Tabellenmodell als Anlage~\ref{anh:diagramme} beigefuegt.]

\section{Projektinitialisierung und Grundkonfiguration}
[PLATZHALTER: Aufbau des Projektgeruests, Konfiguration, Basiskomponenten.]

\section{[PLATZHALTER: Kernmodul 1 -- Name einsetzen]}
[PLATZHALTER: Umsetzung des ersten fachlichen Kernbausteins.]

\section{[PLATZHALTER: Kernmodul 2 -- Name einsetzen]}
[PLATZHALTER: Umsetzung des zweiten fachlichen Kernbausteins.]

\section{Benutzeroberflaeche}
[PLATZHALTER: Umsetzung der UI/UX, Bezug zu den Mockups aus
Anlage~\ref{anh:mockups}. Screenshots der fertigen Anwendung als
Anlage~\ref{anh:screenshots} beigefuegt.]

\section{Hilfsfunktionen und weitere Logik}
[PLATZHALTER: Uebergreifende Hilfsfunktionen, Fehlerbehandlung, Logging o.Ae.]

\chapter{Testphase und Fehlerbehebung}
\label{ch:testphase_und_fehlerbehebung}

[PLATZHALTER: Durchfuehrung der in
Abschnitt~\ref{ch:vorgehensmodell} definierten Testfaelle; gefundene Fehler
und deren Behebung nachvollziehbar darstellen (nicht nur "wurde getestet",
sondern welche Fehler mit welcher Ursache und welcher Loesung). Vollstaendiges
Testprotokoll als Anlage~\ref{anh:testprotokoll} beigefuegt.]

\chapter{Bereitstellung und Abnahme}
\label{ch:bereitstellung_und_abnahme}

\section{Bereitstellung und Inbetriebnahme}
[PLATZHALTER: Deployment/Rollout, Uebergabe an den Fachbereich/Kunden,
Einweisung/Schulung.]

\section{Soll-Ist-Vergleich}
\label{sec:soll_ist_vergleich}
% Vorgabe: "Abweichungen und Anpassungen vom Projektantrag sollen
% nachvollziehbar dargestellt und begruendet werden."
[PLATZHALTER: Abgleich der tatsaechlichen Umsetzung (Zeit, Umfang, Ergebnis)
mit dem genehmigten Projektantrag (Anlage~\ref{anh:projektantrag}). Jede
Abweichung einzeln benennen und begruenden -- nicht nur pauschal
feststellen, dass abgewichen wurde.]

\section{Interne Abnahme des Projekts und Dokumentation}
[PLATZHALTER: Formale Abnahme durch Fachbereich/Ausbilder/in; Hinweis auf
beigefuegte Entwickler- (Anlage~\ref{anh:entwicklerdokumentation}) und
Kundendokumentation (Anlage~\ref{anh:kundendokumentation}).]

\chapter{Abschluss}
\label{ch:abschluss}

\section{Retrospektive und Lessons Learned}
[PLATZHALTER: Was ist gut gelaufen, was wuerde man anders machen? Bezug zu
den in Kapitel~\ref{ch:projektplanung} identifizierten Risiken herstellen:
sind sie eingetreten, wie wurde reagiert?]

\section{Ausblick}
[PLATZHALTER: Moegliche Weiterentwicklungen, offene Punkte, die bewusst
ausserhalb der Projektabgrenzung (Abschnitt~\ref{ch:einleitung}) geblieben
sind.]


% =============================================================================
% NACHTEXT (roemische Seitenzahlen, Fortsetzung der Verzeichnis-Zaehlung)
% =============================================================================
\pagenumbering{Roman}
\setcounter{page}{\value{romanWeiter}}

\chapter*{Literaturverzeichnis}
\addcontentsline{toc}{chapter}{Literaturverzeichnis}

\begin{thebibliography}{99}

\bibitem{platzhalter1}
[PLATZHALTER: Autor/Organisation] ([PLATZHALTER: Jahr]):
\emph{[PLATZHALTER: Titel]}.
\url{[PLATZHALTER: URL]},
Zugriff am [PLATZHALTER: TT.MM.JJJJ].

\end{thebibliography}

\chapter*{Anhang}
\addcontentsline{toc}{chapter}{Anhang}

\section*{Hinweis zu den Anlagen}
Die folgenden Anlagen sind gem\"a\ss{} Vorgabe beizuf\"ugen. Regeln:
\begin{itemize}
  \item Gesamtumfang aller Anlagen: max.\ 20 Seiten.
  \item Auf jede hier gelistete Anlage muss im Hauptteil eindeutig verwiesen
    werden (per \verb|\ref{}|, siehe Beispiele in den Kapiteln).
  \item Anlagen, die nicht selbst erstellt wurden (offizielle Formulare
    o.\,\"a.), sind als solche zu kennzeichnen -- siehe Anlage~1 und Anlage~18.
  \item Vertrauliche/datenschutzrelevante Inhalte in Anlagen sind ebenso zu
    kennzeichnen wie im Hauptteil.
  \item Nicht ben\"otigte Anlagen (z.\,B. wenn im eigenen Projekt kein
    Gantt-Diagramm separat von der Zeitplanung existiert) l\"oschen und die
    zugeh\"orige \verb|\input|-Zeile unten entfernen -- nicht leer stehen
    lassen.
  \item Vor dem Hochladen pr\"ufen, ob eingebundene PDFs (insbesondere
    gescannte Unterschriften) komprimiert sind (Gesamtgr\"o\ss{}e beachten).
\end{itemize}

\section*{Anlage 1: Projektantrag}
\label{anh:projektantrag}
Der urspr\"unglich bei der IHK eingereichte und genehmigte Projektantrag.
Dient als Nachweis, dass die Durchf\"uhrung erst nach Genehmigung begonnen
wurde, und als Referenz f\"ur den Soll-Ist-Vergleich in
Abschnitt~\ref{sec:soll_ist_vergleich}.

[PLATZHALTER: Projektantrag als PDF einbinden, z.\,B.:]
\begin{verbatim}
% \includepdf[pages=-]{anlagen/projektantrag.pdf}
\end{verbatim}

\section*{Anlage 2: Lastenheft}
\label{anh:lastenheft}
Referenziert in Abschnitt~\ref{sec:lastenheft}.

[PLATZHALTER: Vollst\"andiges Lastenheft (Anforderungen aus Nutzersicht)
einf\"ugen -- entweder als eigenes Kapitel hier oder als eingebundenes PDF.]

\section*{Anlage 3: Pflichtenheft}
\label{anh:pflichtenheft}
Referenziert in Abschnitt~\ref{sec:pflichtenheft}.

[PLATZHALTER: Vollst\"andiges Pflichtenheft (technische Umsetzungsvorgaben)
einf\"ugen.]

\section*{Anlage 4: Gantt-Diagramm}
\label{anh:gantt}
Referenziert in Abschnitt~\ref{sec:zeitplanung}.

[PLATZHALTER: Grafisches Gantt-Diagramm der Projektphasen einf\"ugen, z.\,B.:]
\begin{verbatim}
% \includegraphics[width=\linewidth]{bilder/gantt.png}
\end{verbatim}

\section*{Anlage 5: Zeitplanung}
\label{anh:zeitplanung}
Referenziert in Abschnitt~\ref{sec:zeitplanung}.

[PLATZHALTER: Detaillierte Zeitplanung (Phasen, geplante/tats\"achliche
Stunden) als Tabelle einf\"ugen.]

\begin{table}[h!]
  \centering
  \caption{Zeitplanung}
  \begin{tabular}{@{}p{0.6\linewidth}rr@{}}
    \toprule
    \textbf{Projektphase} & \textbf{Geplant (h)} & \textbf{Ist (h)} \\
    \midrule
    {[}PLATZHALTER{]} & {[}PLATZHALTER{]} & {[}PLATZHALTER{]} \\
    \bottomrule
  \end{tabular}
\end{table}

\section*{Anlage 6: Kostenplanung}
\label{anh:kostenplanung}
Referenziert in Abschnitt~\ref{ch:projektplanung}.

[PLATZHALTER: Kostenaufstellung (Personal-, Sach-, Lizenzkosten) und ggf.
Amortisationsrechnung/Wirtschaftlichkeitsbetrachtung einf\"ugen.]

\section*{Anlage 7: Risikoanalyse}
\label{anh:risikoanalyse}
Referenziert in Abschnitt~\ref{ch:projektplanung}.

[PLATZHALTER: Vollst\"andige Risikoanalyse (Risiko, Eintrittswahrscheinlichkeit,
Auswirkung, Gegenma\ss{}nahme) als Tabelle einf\"ugen.]

\section*{Anlage 8: Stakeholderanalyse}
\label{anh:stakeholderanalyse}
Referenziert in Abschnitt~\ref{ch:projektplanung}.

[PLATZHALTER: Stakeholdermatrix (Einfluss/Interesse) einf\"ugen.]

\section*{Anlage 9: Vorgehensmodell}
\label{anh:vorgehensmodell}
Referenziert in Abschnitt~\ref{ch:vorgehensmodell}.

[PLATZHALTER: Grafische Darstellung des gew\"ahlten Vorgehensmodells
(z.\,B. Phasenmodell, Sprint-\"Ubersicht) einf\"ugen.]

\section*{Anlage 10: Testprotokoll}
\label{anh:testprotokoll}
Referenziert in Kapitel~\ref{ch:testphase_und_fehlerbehebung}.

[PLATZHALTER: Vollst\"andiges Testprotokoll (Testfall, erwartetes Ergebnis,
tats\"achliches Ergebnis, Status) als Tabelle einf\"ugen.]

\section*{Anlage 11: Diagramme (Use-Case, ER-/Tabellenmodell, Aktivit\"atsdiagramme)}
\label{anh:diagramme}
Referenziert in Kapitel~\ref{ch:umsetzung}.

[PLATZHALTER: Use-Case-Diagramm, ER-Modell/Tabellenmodell,
Komponenten-/Klassendiagramm sowie Aktivit\"atsdiagramme der wichtigsten
Abl\"aufe einf\"ugen (je Diagrammtyp eigene \verb|\subsection*|, falls mehr
als eines).]

\section*{Anlage 12: Mockups}
\label{anh:mockups}
Referenziert in Abschnitt~\ref{ch:umsetzung}.

[PLATZHALTER: UI-Entw\"urfe/Mockups vor der Umsetzung einf\"ugen.]

\section*{Anlage 13: Screenshots der Anwendung}
\label{anh:screenshots}
Referenziert in Abschnitt~\ref{ch:umsetzung}.

[PLATZHALTER: Screenshots der fertigen Anwendung einf\"ugen. Bei
Kundendaten/internen Daten in Screenshots: unkenntlich machen oder als
vertraulich kennzeichnen.]

\section*{Anlage 14: Entwicklerdokumentation}
\label{anh:entwicklerdokumentation}
Referenziert in Abschnitt~\ref{sec:soll_ist_vergleich}.

[PLATZHALTER: Technische Dokumentation f\"ur Nachfolger/Wartung
(Installation, Architektur, wichtige Klassen/Module) einf\"ugen.]

\section*{Anlage 15: Kundendokumentation}
\label{anh:kundendokumentation}
Referenziert in Abschnitt~\ref{sec:soll_ist_vergleich}.

[PLATZHALTER: Anleitung f\"ur Endnutzer (Bedienung der Anwendung) einf\"ugen.]

\section*{Anlage 16: Quellcode-Ausschnitte}
\label{anh:quellcode}
Referenziert in Kapitel~\ref{ch:umsetzung}.

[PLATZHALTER: Ausgew\"ahlte, besonders erkl\"arungsbed\"urftige
Quelltext-Ausschnitte einf\"ugen -- nicht der vollst\"andige Quelltext
(Seitenlimit). Z.\,B.:]
\begin{verbatim}
% \lstinputlisting[language=Python, firstline=1, lastline=40]{code/beispiel.py}
\end{verbatim}

\section*{Anlage 17: Protokoll zur Projektarbeit}
\label{anh:protokoll_projektarbeit}
Nicht selbst erstellt -- offizielles, von der IHK bereitgestelltes Formular
(zum Herunterladen auf der IHK-Downloadseite verf\"ugbar). Als
unterschriebenes PDF einbinden:
\begin{verbatim}
% \includepdf[pages=-]{anlagen/protokoll_projektarbeit.pdf}
\end{verbatim}
[PLATZHALTER: Datei ablegen unter \texttt{anlagen/protokoll\_projektarbeit.pdf}
und obige Zeile aktivieren.]

\section*{Anlage 18: Pers\"onliche Erkl\"arung}
\label{anh:persoenliche_erklaerung}
Nicht selbst erstellt -- offizielles, von der IHK bereitgestelltes Formular
(zum Herunterladen auf der IHK-Downloadseite verf\"ugbar). Als
unterschriebenes PDF einbinden:
\begin{verbatim}
% \includepdf[pages=-]{anlagen/persoenliche_erklaerung.pdf}
\end{verbatim}
[PLATZHALTER: Datei ablegen unter \texttt{anlagen/persoenliche\_erklaerung.pdf}
und obige Zeile aktivieren.]



\end{document}
